% Open questions for the author (ma.tex):
% 1. Lemma ma follows Tan's Result 2.7 (arXiv:1610.01914v4) as quoted in R20 ("every nontrivial
%    character"). No review opened Pott's book (R20: "no text available"), so the section number
%    \S1.5 comes from solutions.md and is unchecked. Ma's 1985 thesis is not in refs.bib; I did not
%    add it, since its title is not recorded anywhere in the repo.
% 2. Tan's numbers (Results 2.6-2.7, Thm 4.22-4.23) are those of arXiv v4 (R23, R28); v1 numbers
%    them differently (R23: 2.13, 2.15).
% 3. Schmidt-Smith Lemma 6.2: R28 could not read the paper (only Tan's use of it), and P1 says
%    "there for weight 36". The remark says only "for weighing matrices"; please check the lemma.
% 4. Theorem z32 says "if no Z exists"; Lean's hypothesis is a kernel-run exhaustive search
%    (noVector) that Lean proves sound, not complete. Lemma lem:halved is proved inside the Lean
%    proofs of Theorems z29b and z32, not as a named Lean theorem. Is this wording acceptable?
% 5. \mathcal{P} is defined in sec:characters (eq:cyclo); see the note there.
% 6. \zrows{29} is a total (solutions.md: 177 by the parity, 94 by the halving).

\section{Ma's lemma}\label{sec:ma}

When $p$ divides $|H|$, Turyn's lemma still gives $p^b \mid \chi(a)$ for every nontrivial
character, provided $p$ is self-conjugate modulo the exponent of $H$. Ma's lemma turns this into
a statement about the coset sums themselves. We identify a vector $a$ on $H$ with
$\sum_h a_h h \in \Z[H]$, and a subgroup $P$ with $\sum_{u \in P} u$.

\begin{lemma}[Ma; see~{\cite[\S1.5]{pott}} and~{\cite[Result~2.7]{tan}}]\label{lem:ma}
Let $p$ be a prime and $H$ a finite abelian group whose Sylow $p$-subgroup is cyclic and
nontrivial, and let $P$ be its subgroup of order $p$. If $Y \in \Z[H]$ and $p^b$ divides $\chi(Y)$
for every nontrivial character $\chi$ of $H$, then $Y = p^b X_1 + P X_2$ for some
$X_1, X_2 \in \Z[H]$.
\end{lemma}

Throughout this section, $p$ is a prime, $q = p^j \in \mathcal{P}$ (see~\eqref{eq:cyclo}), and
$K$ is a finite abelian group with $p^i \equiv -1 \pmod{\exp K}$ for some $i \ge 1$ ($K = 0$ is
allowed). Then $p \nmid |K|$, and $p$ is self-conjugate modulo the exponent of
$H = C_q \times K$. Let $\phi \colon G \to H$ be onto, $c = v/|H|$, $u = (q/p, 0)$, and $P$ the
subgroup of order $p$, generated by $u$.

\begin{lemma}\label{lem:mastep}
Let $b \ge 0$ with $p^{2b} \mid n$. If $a$ is the vector of coset sums of an SDS on $H$, then
$p^b$ divides $a_{h+u} - a_h$ for every $h \in H$.
\end{lemma}

\begin{proof}
For $\chi \ne \chi_0$ the value $\chi(a)$ lies in $\Z[\zeta_{\exp H}]$ and
$\chi(a)\overline{\chi(a)} = n$ (Lemma~\ref{lem:chars}), so $p^b \mid \chi(a)$ by
Lemma~\ref{lem:turyn}. The Sylow $p$-subgroup $C_q$ of $H$ is cyclic, so Lemma~\ref{lem:ma} gives
$a = p^b X_1 + P X_2$. The coefficients of $P X_2$ are constant on the cosets of $P$, hence
$a_{h+u} - a_h = p^b \bigl(X_1(h+u) - X_1(h)\bigr)$.
\end{proof}

The Lean proof (\lean{SignedDifferenceSets.Ma.dvd_contract_sub}) uses neither lemma as a black
box; it proves the case at hand directly. Let $A = \Z[K][X]/(\Phi_q(X))$ and
$\psi \colon \Z[H] \to A$ the ring homomorphism with $(x, \kappa) \mapsto X^x \kappa$; it is onto,
and every element of its kernel is constant on the cosets of $P$. Since $\psi(H) = 0$, the
identity $a\,a^{(-1)} = n + \lambda c H$ gives $\psi(a)\,\sigma(\psi(a)) = n$, where $\sigma$ is the
involution $X \mapsto X^{-1}$, $\kappa \mapsto -\kappa$. The element $\pi = X - 1$ is a
non-zero-divisor, $A/\pi A \cong \mathbb{F}_p[K]$, and there $\sigma$ is $\kappa \mapsto -\kappa =
p^i\kappa$, that is, $z \mapsto z^{p^i}$. As $\mathbb{F}_p[K]$ is reduced, $z\,\sigma(z) \in \pi A$
forces $z \in \pi A$; since $\sigma(\pi) = -X^{-1}\pi$, induction gives $\psi(a) \in
\pi^{b\varphi(q)} A = p^b A$, the last step by the cyclotomic input of
Section~\ref{sec:sylow}. So $a - p^b X_1 \in \ker\psi$ for some $X_1$, which is the conclusion of
Lemma~\ref{lem:ma} for this $a$.

\begin{theorem}\label{thm:z24}
Suppose $n \ne 0$, $p^{2b} \mid n$ and $p^b > 2c$. Then $G$ has no SDS with parameters
$(v,k,\lambda)$.
\end{theorem}

\begin{proof}
Let $a$ be the coset sums of an SDS on $H$. By Lemma~\ref{lem:mastep}, $p^b$ divides
$a_{h+u} - a_h$, whose absolute value is at most $2c < p^b$ (Lemma~\ref{lem:contraction}). So $a$
is constant on the cosets of $P$. A character $\chi$ with $\chi(u) \ne 1$ then has
$\chi(a) = 0$, since $\sum_{l < p} \chi(u)^l = 0$; this contradicts $|\chi(a)|^2 = n \ne 0$.
\end{proof}

\begin{remark}\label{rem:z24}
The inequality $p^b > 2c$ cannot be relaxed to $p^b \ge 2c$: the $(4,4,0)$ set $(1,1,1,-1)$ in
$C_4$ has $p^b = 2 = 2c$. Example: $(27,17,8)$ in $C_{27}$ (N75) has $n = 9$ and
$s^2 = 17 + 8 \cdot 26 = 225$. With $H = G$ ($q = 27$, $K = 0$, $c = 1$) and $b = 1$ we have
$3 > 2$: by Lemma~\ref{lem:mastep}, $D(x+9) \equiv D(x) \pmod 3$, so $D$ would have period $9$,
and every character of order $27$ would vanish on $D$.

Origin: Ma's lemma (Ma 1985; see~\cite{pott}) with Turyn's self-conjugacy~\cite{turyn}, a
combination known for circulant weighing matrices~\cite[Results~2.6--2.7]{tan}; the form for SDS
on quotients $C_q \times K$ and its Lean proof are ours. It settles \zrows{24} cases (Lean: the
statement \lean{SignedDifferenceSets.MaRule}, proved as \lean{SignedDifferenceSets.maRule}).
\end{remark}

\subsection{The edge case $p = 2$, $b = 1$}

Now let $p = 2$, so $q = 2r$ and $|K|$ is odd, and suppose $4 \mid n$. Let $L = C_r \times K$ and
$\psi \colon H \to L$ the reduction, with kernel $P = \{0, u\}$. Lemma~\ref{lem:mastep} with
$b = 1$ says that $a_h$ and $a_{h+u}$ have the same parity, which Theorem~\ref{thm:z24} cannot use
($2 > 2c$ fails). Adding them gives the following.

\begin{lemma}\label{lem:halved}
Let $D$ be an SDS with parameters $(v,k,\lambda)$ on $G$ and $W$ its coset sums on $L$ (along
$\psi\phi$, with fibres of size $2c$). Then every $W_l$ is even, and $Z = W/2$ satisfies
\begin{enumerate}
\item $|Z_l| \le c$ for every $l$;
\item $\sum_l Z_l^2 = (k + \lambda(2c - 1))/4$;
\item $\sum_l Z_l Z_{l+g} = \lambda c/2$ for every $g \ne 0$ in $L$;
\item $\sum_l Z_l = \pm s/2$;
\item $\sum_l |Z_l| \le k/2$ and $\sum_l |Z_l| \equiv s/2 \pmod 2$.
\end{enumerate}
\end{lemma}

\begin{proof}
For $h \in H$, $W_{\psi(h)} = a_h + a_{h+u} = 2a_h + (a_{h+u} - a_h)$ is even by
Lemma~\ref{lem:mastep}. Items (1)--(4) are Lemma~\ref{lem:contraction} for $L$, whose fibres have
size $2c$, divided by $2$, $4$, $4$ and $2$. For (5), $\sum_l |W_l| \le \sum_x |D(x)| = k$, and
$|Z_l| \equiv Z_l \pmod 2$.
\end{proof}

\begin{theorem}\label{thm:z29a}
If $G$ has an SDS with parameters $(v,k,\lambda)$, then $8$ divides
$2s - (k + \lambda(2c-1))$ or $-2s - (k + \lambda(2c-1))$.
\end{theorem}

\begin{proof}
By Lemma~\ref{lem:halved}(2), (4) and $z^2 \equiv z \pmod 2$,
$(k + \lambda(2c-1))/4 = \sum_l Z_l^2 \equiv \sum_l Z_l = s'/2 \pmod 2$ for some
$s' \in \{s, -s\}$; multiply by $4$.
\end{proof}

The Lean proof sums $2F - F^2$ over $H$, with $F_h = a_h + a_{h+u}$, and pairs $h$ with $h + u$;
this is the same computation.

\begin{theorem}\label{thm:z29b}
Suppose $c = 1$, so that $\phi$ is an isomorphism $G \cong C_{2r} \times K$, and let
$k' = (k + \lambda)/4$ and $\lambda' = \lambda/2$ be integers, $k' \ge 0$. If $G$ has an SDS with parameters
$(v,k,\lambda)$, then $C_r \times K$ has an SDS with parameters $(v/2, k', \lambda')$.
\end{theorem}

\begin{proof}
For $c = 1$, Lemma~\ref{lem:halved}(1) gives $Z_l \in \{0, \pm 1\}$, (2) says that $Z$ has $k'$
nonzero values, and (3) that its autocorrelation is $\lambda'$.
\end{proof}

\begin{remark}\label{rem:z29}
For $c = 1$, Theorem~\ref{thm:z29a} is the parity condition $k' \equiv s/2 \pmod 2$ of the halved
set. The cases settled by Theorem~\ref{thm:z29b} are those whose halved case has no SDS by a Lean
proof (\lean{SignedDifferenceSets.HalfCert.not_hasSDS}). Examples: $(48,42,6)$ in $C_{48}$ (N343)
has $n = 36$ and $s^2 = 42 + 6 \cdot 47 = 324$, $s = 18$. On $H = C_{16}$ ($K = 0$, $c = 3$),
$k + \lambda(2c-1) = 72$, and $\pm 36 - 72 \equiv 4 \pmod 8$, so Theorem~\ref{thm:z29a} excludes
it. And $(72,40,4)$ in $C_{72} \cong C_8 \times C_9$ (N935) has $n = 36$; with $K = C_9$ ($2^3
\equiv -1 \pmod 9$) and $c = 1$, Theorem~\ref{thm:z29b} gives a $(36,11,2)$ set in $C_4 \times
C_9 \cong C_{36}$ (case N171). There is none, by Theorem~\ref{thm:z24} with $p = 3$, $H = C_9
\times C_4$ ($3 \equiv -1 \pmod 4$), $c = 1$, $n = 9$ and $b = 1$.

Origin: Ma's lemma (Ma 1985). For weighing matrices the halving at $c = 1$ is in Schmidt and
Smith~\cite[Lemma~6.2]{schmidt-smith}, as used in~\cite[Thms.~4.22--4.23]{tan}; the parity for
general $c$ is ours; we did not find it in print. Together the two theorems settle \zrows{29}
cases (Lean: \lean{SignedDifferenceSets.not_hasSDS_of_maDescent} and
\lean{SignedDifferenceSets.hasSDS_half_of_maDescent}).
\end{remark}

\begin{theorem}\label{thm:z32}
Suppose that $s/2$, $(k + \lambda(2c-1))/4$ and $\lambda c/2$ are integers and that no function
$Z \colon L \to \Z$ satisfies (1)--(5) of Lemma~\ref{lem:halved}. Then $G$ has no SDS with
parameters $(v,k,\lambda)$.
\end{theorem}

\begin{proof}
Otherwise Lemma~\ref{lem:halved} provides such a $Z$.
\end{proof}

\begin{remark}\label{rem:z32}
In Lean (\lean{SignedDifferenceSets.not_hasSDS_of_evenContraction}) the hypothesis that no $Z$
exists is the result of an exhaustive search, which Lean proves sound and runs in the kernel for
each case. Example: $(64,18,2)$ in $C_2 \times C_{32}$ (N664) has $n = 16$ and $s^2 = 18 + 2 \cdot
63 = 144$, $s = 12$. Take $H = C_{32}$ ($K = 0$, $c = 2$) and $L = C_{16}$. Then $Z$ would have
$|Z_l| \le 2$, $\sum_l Z_l^2 = 6$, autocorrelation $2$ and $\sum_l Z_l = \pm 6$. Since $|z| \le
z^2$ for integers, $6 \le \sum_l |Z_l| \le \sum_l Z_l^2 = 6$, so $\pm Z$ is the indicator of a
$(16,6,2)$ difference set in $C_{16}$. None exists: such a set $E$ has $|\chi(E)|^2 = 4$ for
$\chi \ne \chi_0$, so $2 \mid \chi(E)$ by Lemma~\ref{lem:turyn}, and Lemma~\ref{lem:ma} gives
$E(x+8) \equiv E(x) \pmod 2$, hence $E(x+8) = E(x)$, and every character of order $16$ would
vanish on $E$. Lean settles this case, like the others, by the search.

Origin: Ma's lemma (Ma 1985) and Turyn's self-conjugacy~\cite{turyn}; the contraction for SDS and
its Lean proof are ours. It settles \zrows{32} cases.
\end{remark}
